% ---------- Capa ---------------------------------------------------------------
\begin{titlepage}
\begin{tikzpicture}[remember picture, overlay]
  \fill[azulMB] (current page.south west) rectangle (current page.north east);
  \begin{scope}
    \clip (current page.north west) rectangle ([yshift=-0.42\paperheight]current page.north east);
    \node[anchor=north, inner sep=0pt] at (current page.north)
      {\includegraphics[height=0.42\paperheight]{\imagens/ilha_das_cobras.jpg}};
  \end{scope}
  \fill[ouroMB] ([yshift=-0.42\paperheight]current page.north west) rectangle ([yshift=-0.42\paperheight-1.6mm]current page.north east);
  \node[anchor=north west, text=ouroMB, font=\bfseries] (k) at ([xshift=22mm,yshift=-0.42\paperheight-16mm]current page.north west)
    {CAPACITAÇÃO INTERNA \,·\, APOSTILA-GUIA};
  \node[anchor=north west, text=white, align=left, inner xsep=0pt,
        font=\titulofonte\fontsize{44}{46}\selectfont] (t) at ([yshift=-3mm]k.south west)
    {GIT, GITLAB E IA APLICADA\\À ENGENHARIA DE SOFTWARE};
  \node[anchor=north west, text=white!82!azulMB, align=left, inner xsep=0pt, text width=15cm,
        font=\montlight\large] (s) at ([yshift=-4mm]t.south west)
    {Do primeiro commit à documentação de sistemas legados com um agente de IA: um guia para quem está começando};
  \node[anchor=south west, text=white, align=left, inner xsep=0pt, font=\normalsize] at ([xshift=22mm,yshift=22mm]current page.south west)
    {\textbf{Ronivaldo D. Andrade}\\[2pt]
     {\color{white!75!azulMB}\small Estagiário de Análise e Desenvolvimento de Sistemas}\\
     {\color{white!75!azulMB}\small Marinha do Brasil}\\[6pt]
     {\color{ouroMB}\small\bfseries Capacitação em 01/10/2026, das 10h às 11h}};
  \node[anchor=south east, inner sep=0pt] at ([xshift=-20mm,yshift=22mm]current page.south east)
    {\marca[white]{30}};
\end{tikzpicture}
\end{titlepage}

% ---------- Sobre esta apostila -------------------------------------------------
\chapter*{Sobre esta apostila}
\markboth{Sobre esta apostila}{}

Esta apostila acompanha a capacitação \textbf{\enquote{Git, GitLab e IA aplicada à Engenharia de Software}}, apresentada na Marinha do Brasil. A apresentação tem uma hora e é só expositiva. Aqui está tudo o que foi dito com mais calma, os detalhes que não cabem num slide e os \textbf{exercícios práticos}, para você aprender no seu ritmo.

Ela foi escrita para quem \textbf{nunca usou Git}. Se você já usa, pule direto para os capítulos \ref{cap:problema1} e \ref{cap:problema2}, que contam os dois problemas reais resolvidos no estágio.

\begin{center}
\renewcommand{\arraystretch}{1.35}
\begin{tabularx}{0.92\linewidth}{@{}>{\bfseries\color{azulMB}}l X@{}}
\toprule
Capacitação & Git, GitLab e IA aplicada à Engenharia de Software \\
Autor e instrutor & Ronivaldo D. Andrade \\
Função & Estagiário de Análise e Desenvolvimento de Sistemas \\
Organização & Marinha do Brasil \\
Data da capacitação & 01/10/2026, das 10h às 11h \\
Data de elaboração & 28/09/2026 \\
Material complementar & Apresentação de slides (\arq{apresentacao/apresentacao.pdf}) \\
\bottomrule
\end{tabularx}
\end{center}

\section*{Como ler esta apostila}

Ao longo do texto aparecem caixas coloridas. Cada cor tem um papel:

\begin{conceito}[Conceito]
Uma definição importante, para guardar.
\end{conceito}
\begin{analogia}
Uma comparação com o dia a dia, para a ideia ficar concreta.
\end{analogia}
\begin{dica}
Um atalho ou uma boa prática.
\end{dica}
\begin{atencao}
Um erro comum ou um cuidado que evita dor de cabeça.
\end{atencao}

\begin{terminal}
# Isto é um comando para digitar no terminal.
# As linhas que começam com # são comentários: não precisa digitar.
git status
\end{terminal}

\begin{saida}
Isto é o que o computador responde. O texto exato pode variar
conforme a versão e o idioma do seu Git.
\end{saida}

Quando um comando traz algo entre \arq{<sinais de menor e maior>}, como \arq{<URL\_DO\_REPOSITORIO>}, troque tudo, inclusive os sinais, pelo valor real.

\begin{atencao}[Antes de usar IA com código da instituição]
O capítulo \ref{cap:problema2} mostra como um agente de IA ajudou a documentar sistemas. Um agente desse tipo envia trechos do código para um serviço externo. \textbf{Só use com sistemas institucionais quando houver autorização e respeitando as normas de segurança da informação da sua OM.} Para praticar, use projetos pessoais ou de código aberto.
\end{atencao}
